\section{Introduction}

Spoken interaction with graphical interfaces combines automatic speech recognition (ASR) with visual grounding: the system must recover a command and locate its intended control. Thai commands may mix Thai syntax with English screen labels, paraphrase a label, describe an intent, or refer to a visual position. A small transcript error may remove a useful screen anchor, while visual context may still support a correct click.

ThaiVoice-GUIBench makes these outcomes inspectable in a fixed single-click task. It contains 20 synthetic mobile-style screens and TTS speech, rather than production interfaces or real-speaker recordings. Four command tiers distinguish verbatim labels, paraphrases, intents, and visual/positional references. We compare two ASR systems and two grounding models, then evaluate three fixed-rule transcript corrections; the release includes the pipeline diagram. Grounding prompts receive the screenshot and command, without the target identifier, target box, gold answer, or candidate list.

Correction is evaluated at both transcript and click levels. Visible screen strings may help recover a label but may also inject unsupported content. Moreover, raw and corrected grounding were run independently, days apart. A changed click can therefore reflect rerun variability as well as changed text. Changed/unchanged subsets and no-op controls reveal this attribution problem; they do not solve it.

We ask how ASR errors vary by tier and TTS condition, whether raw-ASR grounding differs detectably from gold-text grounding, and when transcript changes coincide with better or worse clicks. Our contributions are an auditable Thai diagnostic resource, tiered evaluation that retains parse and out-of-bounds failures, three frozen correction baselines with benefit-and-harm reporting, and analyses exposing rerun confounding. The results support diagnosis within this synthetic benchmark, without establishing effects for real users or broader code-switching patterns.
